\section{Threats to Validity}
\label{sec:threats}

\subsection{Sample Size and Ties}

The unit of analysis is the service, so the sample sizes are \val{power_sn_n} for Social Network
and \val{power_hr_n} for Hotel Reservation, the number of services the applications have. Both
were registered in advance and matched exactly. %[id:power_sn_n_matches] %[id:power_hr_n_matches]
Repetition makes each service's outcome more precise, which is why each cell was run
\val{campaign_reps_per_cell} times, but it adds no points to a correlation. Ties tighten the
bound further: the caller count takes \val{power_sn_levels} distinct values on Social Network
and \val{power_hr_levels} on Hotel Reservation, which caps the correlation for an outcome
without ties at \val{power_sn_ceiling} and \val{power_hr_ceiling}. These ceilings were stated
before data collection, so they are a property of the design and not an explanation found
afterwards. On Hotel Reservation the outcome was also fully saturated, so every result there
restates how the outcome was built and is reported as directional only.
%[id:saturation_hr_kill_tautological]
Its value is as a second, differently shaped topology on which the direction result also holds.

\subsection{A Predictor Close to Its Outcome}

The most serious threat to the primary result is conceptual. The outcome counts degraded
transitive callers and the predictor counts transitive callers, so they coincide whenever every
ancestor degrades, and on Social Network the whole gap between $\rho = \val{primary_rho}$ and a
correlation of one rests on user-service. Withholding it raises $\rho$ to
\val{loo_userservice_rho}, a change \val{loo_swing_ratio} times the size of the next largest,
in a check specified after the result was known. The saturation measure was registered before
collection precisely to detect this, and the test remained meaningful because saturation was
not guaranteed in advance and did not fully hold. A reader may still treat the Social Network
correlation as one service's worth of evidence.

\subsection{Benchmarks, Workload and Graph}
\label{sec:threat-scope}

Both applications come from one benchmark suite at one commit, chosen because they are
instrumented and reproducible, not because they sample the space of microservice architectures.
The workload is synthetic, fixed at \val{protocol_load_users} users with a fixed request mix,
and containment mechanisms were not studied. Nothing here describes production behaviour.
Because the dependency graph is rebuilt from observed traffic, it is a property of deployment and
workload together; production call graphs change from request to request \cite{luo2021}. Even
here, the traces recover \val{hr_manifest_edges_recovered} of the
\val{hr_manifest_expected_edges} edges implied by Hotel Reservation's documentation, and the user
service has one observed caller where the documentation implies two.
%[id:hr_user_observed_caller_count]
A different request mix would give it a different predictor value with nothing about the
deployed system changed, which limits how far any relationship reported here carries to other
traffic.

\subsection{Measurement}

For latency faults, the runner infers that the netem rule is gone from the exit of the injecting
process. That inference failed at least once (Deviation 4). The cross-run baseline gate added
afterwards catches a rule that outlives its run, because later baselines then look contaminated,
but a rule that leaked and cleared within a single run would pass unnoticed, so latency-fault
recovery times carry an assumption the instrumentation cannot check. Blast-radius counts are
less exposed, since they come from the fault window itself. The server-span replay of
Deviation 11 could not cover recovery times on Social Network, because the probe's windows were
not stored, and those times stand as measured. The baseline window includes the load ramp-up (Deviation 13); the same
procedure ran in every run and the baseline gate reported no exceptions,
%[id:integrity_checks_passed]
but the size of any start-up effect was not measured. Two fault types were studied, a process
kill and a \val{fault_latency_delay_ms}~ms delay; CPU contention was evaluated in the pilot and
dropped because its effect was too small to classify reliably. Finally, some testbed values in
Table~\ref{tab:tbl_ieee_testbed} were queried after the campaign, on \val{testbed_measured_on},
and no record from the campaign days confirms that the host operating system build and Docker
Desktop version were unchanged.
